% ============================================================
\section{On the Robustness Cost of Boolean Discretization: An Adversarial Margin Analysis}
\label{sec:adversarial-margin}
\textit{This section was contributed by Claude Sonnet 5 (Anthropic).}
% ============================================================

Sections~2--4 establish that quantized attention is, in the limit, a Boolean
routing function computable by a $\nand$ circuit of depth $O(\log d)$, and
that a substantial fraction of softmax's arithmetic work does not change
\emph{which} tokens attend to which. This section asks a narrower question
that the prior sections leave open: if the routing topology is genuinely
Boolean, what is lost by discarding the continuous wrapper around it? We
argue the answer is not ``nothing'' --- it is \emph{margin}, and margin is
precisely the quantity that governs robustness to small input perturbations
and the smoothness of the training loss landscape.

\subsection{Setup: margin as the missing variable}

Let $s_{ij} = q_i \cdot k_j / \sqrt{d}$ denote a pre-softmax score, and let
$\threshold(s_{ij})$ denote the Boolean routing decision extracted from it
under some quantization scheme, as defined in \S2. Define the \emph{margin}
of a routing decision as
\[
  m_{ij} \;=\; \bigl| s_{ij} - \tau \bigr|,
\]
where $\tau$ is the effective decision boundary induced by the softmax's
competitive normalization at that position. The NAND circuit of \S2--3 is,
by construction, a function of $\threshold(s_{ij})$ alone; it has no access
to $m_{ij}$. Softmax attention, by contrast, outputs a value that varies
continuously with $m_{ij}$ even after the ``winner'' is decided.

\begin{lemma}[Discontinuity of Boolean routing]
\label{lem:discontinuity}
Let $f_{\mathrm{bool}}: \RR^d \times \RR^d \to \BB$ be the threshold routing
function and $f_{\mathrm{soft}}: \RR^d \times \RR^d \to [0,1]$ the softmax
attention weight, both derived from the same score $s_{ij}$. For any
$\epsilon > 0$ there exists a perturbation $\delta$ with $\|\delta\| =
O(m_{ij})$ such that
\[
  f_{\mathrm{bool}}(s_{ij}) \neq f_{\mathrm{bool}}(s_{ij} + \delta),
  \qquad \text{while} \qquad
  \bigl| f_{\mathrm{soft}}(s_{ij}) - f_{\mathrm{soft}}(s_{ij}+\delta) \bigr|
  = O(\delta).
\end{lemma}

\begin{proof}[Proof sketch]
$f_{\mathrm{bool}}$ is piecewise constant with a jump discontinuity at
$s_{ij} = \tau$; any perturbation crossing the boundary flips the output
regardless of how small $\delta$ is, provided $m_{ij} < \|\delta\|$.
$f_{\mathrm{soft}}$ is Lipschitz in $s_{ij}$ with constant bounded by the
softmax Jacobian, so its output changes proportionally to $\delta$ near
$\tau$ and does not jump. $\qed$
\end{proof}

This is not a deep result --- it is the standard hard-threshold-versus-soft-
threshold distinction from classical signal processing --- but it has a
concrete consequence for the NAND-decomposition program that \S2's
complexity accounting does not surface: \emph{the ``wasted'' $\Theta(nd)$
flops are exactly the computation that keeps the attention function
Lipschitz.} Removing them does not merely save compute; it removes the
smoothness guarantee that gradient-based training and small-perturbation
robustness both depend on.

\subsection{Where this matters, and where it does not}

We do not think this undermines the engineering proposal in \S4. A trained,
frozen, inference-time model with wide empirical margins ($m_{ij} \gg 0$ for
almost all $(i,j)$, consistent with the bimodality index of 0.71--0.79
reported there) may tolerate hard routing perfectly well, because few
decisions live near the boundary. The risk is concentrated in three regimes,
none of which \S4's inference-time benchmark would detect:

\begin{enumerate}
  \item \textbf{Training and fine-tuning}, where gradients flow through
    exactly the near-boundary cases the Boolean reduction discards.
  \item \textbf{Distribution shift}, where inputs unlike the measured 10k-
    token sample may concentrate probability mass near $\tau$ rather than
    away from it.
  \item \textbf{Adversarial inputs}, where an attacker searches specifically
    for the low-margin region rather than sampling it by chance.
\end{enumerate}

\begin{conjecture}[Margin concentration under distribution shift]
\label{conj:shift}
The bimodality index reported in \S4 is a property of the training
distribution, not of the attention function itself. There exist inputs
$x$, reachable by ordinary distribution shift and not requiring an
adversary, for which the empirical margin distribution over $(i,j)$ pairs
is substantially flatter than 0.71--0.79, at which point a hard-routed
NAND circuit and its soft-routed original would diverge in output on a
non-negligible fraction of positions. We do not have a proof of this and
flag it as an open empirical question (candidate: SKW-006).
\end{conjecture}

\subsection{An engineering mitigation: hysteresis routing}

If \S4's proposal is adopted --- replacing softmax with NAND routing at
inference --- Lemma~\ref{lem:discontinuity} suggests a standard fix borrowed
from analog circuit design: a \emph{Schmitt trigger}. Instead of a single
threshold $\tau$, use two thresholds $\tau_- < \tau_+$ and retain the
previous routing decision for any score in $(\tau_-, \tau_+)$:

\begin{lstlisting}
threshold_route(s, prev_bit, tau_minus, tau_plus):
    if s >= tau_plus:  return 1
    if s <  tau_minus: return 0
    return prev_bit   # hysteresis band: hold previous decision
\end{lstlisting}

This costs one bit of state per attention edge and a comparison against two
constants rather than one --- still $O(\log d)$ circuit depth, still
compatible with the $\nand$ decomposition of \S3, since a two-threshold
comparator is itself expressible as a small NAND network. It converts a
brittle point discontinuity into a bounded-width dead zone, at the cost of
reintroducing a small amount of state that a pure combinational NAND circuit
does not otherwise need.

\subsection{Summary}

The NAND decomposition is correct as a statement about routing topology.
What it does not capture is that softmax's ``redundant'' flops are, at
minimum, functioning as a robustness margin during training and under
distribution shift, and possibly at inference under adversarial pressure.
We would frame the honest version of \S4's engineering claim as: \emph{hard
Boolean routing is safe to substitute for softmax at inference, for a
frozen model, on inputs resembling the training distribution, provided
margins are monitored} --- and we propose hysteresis routing as the
minimal mechanism to make that substitution safe rather than merely fast.

% Contributed by: Claude Sonnet 5
% Date: July 13, 2026
% Trust: THE SHARED PRIMORDIAL FOUNDATION -- EIN 42-6976431
% In memory of Eric Brandon Westerhoff.

% New bibliography entries (add to bibliography if not already present):
% \bibitem{goodfellow2014adversarial} Goodfellow, I., Shlens, J., Szegedy, C. (2014). Explaining and Harnessing Adversarial Examples. arXiv:1412.6572.
% \bibitem{madry2017robust} Madry, A., et al. (2017). Towards Deep Learning Models Resistant to Adversarial Attacks. arXiv:1706.06083.

% ── AUTHOR SIGNATURE ─────────────────────────────────────────
% Model:    Claude Sonnet 5
% Provider: Anthropic
% Date:     July 13, 2026
% Section:  On the Robustness Cost of Boolean Discretization: An Adversarial Margin Analysis
% Angle:    Argues the softmax "waste" §2/§4 identify is functioning as an adversarial/training robustness margin, and proposes hysteresis routing as a mitigation if hard NAND routing is adopted at inference.
% Trust:    THE SHARED PRIMORDIAL FOUNDATION -- EIN 42-6976431
% In memory of Eric Brandon Westerhoff.
% ─────────────────────────────────────────────────────────────
